% 06-references.tex - References

\chapter{References}
\label{sec:references}

\section{LWA Documentation}

\begin{enumerate}
    \item MCS Common ICD (LWA Engineering Memo MCS0005)
    \item MCS-NDP Interface Control Document
    \item J.~Craig, ``Long Wavelength Array Station Architecture,'' Ver.~2.0, LWA Memo~161, Feb.~26, 2009.
    \item C.~Wolfe, S.~Ellingson, C.~Patterson, ``MCS Data Recorder Preliminary Design \& Verification,'' Ver.~1.0, LWA Memo~165, Aug.~26, 2009.
    \item C.~Wolfe, S.~Ellingson, C.~Patterson, ``MCS-DR Storage Unit,'' LWA Engineering Memo MCS0019, Sep.~23, 2009.
\end{enumerate}

\section{Standards}

\begin{enumerate}
    \item IEC~60027-2, ``Letter symbols to be used in electrical technology --- Part~2: Telecommunications and electronics,'' Third Ed., 2005.
\end{enumerate}

\section*{Software Components}

\begin{longtable}{l p{9cm}}
\toprule
\textbf{Component} & \textbf{Description} \\
\midrule
\endfirsthead
\endhead
\bottomrule
\endfoot

DROS2          & Main Data Recorder Operating Software.  Provides all
                 ICD-specified functionality. \\

Msender        & Command-line utility to send MCS messages to and receive
                 responses from the DR.  Intended for operating the DR in
                 the absence of MCS. \\

DataSource2    & Simulated NDP data source for testing.  Emulates NDP
                 output in DRX, DRX8, TBT, TBS, and COR formats. \\

SpectrogramViewer & GUI utility for viewing spectrometer output files
                    produced by DROS2. \\

build\_helper.py & Build configuration utility for multi-DR deployments.
                   Creates per-instance build directories with separate
                   installation paths and service files. \\

\end{longtable}

\section*{Configuration Files}

\begin{longtable}{l p{9cm}}
\toprule
\textbf{File} & \textbf{Description} \\
\midrule
\endfirsthead
\endhead
\bottomrule
\endfoot

\texttt{defaults\_v2.cfg}
    & Main system configuration (Section~\ref{sec:configuration}). \\

\texttt{binary-mode}
    & Selects the active DROS2 operating mode
      (\texttt{Spectrometer} or \texttt{LiveBuffer}). \\

\texttt{dros.service}
    & \texttt{systemd} unit file for single-DR deployment. \\

\texttt{dros-dr\textit{N}.service}
    & \texttt{systemd} unit file for multi-DR deployment
      (one per DR instance). \\

\end{longtable}

\section{Configuration Parameters}
\label{sec:configuration}

DROS2 configuration is stored in a key-value text file
(\texttt{defaults\_v2.cfg}).  Table~\ref{tab:config} lists the supported
parameters in the order they appear in the configuration file.

\begin{longtable}{l l p{6.5cm}}
\caption{DR Configuration Parameters (\texttt{defaults\_v2.cfg})}
\label{tab:config} \\
\toprule
\textbf{Parameter} & \textbf{Type} & \textbf{Description} \\
\midrule
\endfirsthead

\toprule
\textbf{Parameter} & \textbf{Type} & \textbf{Description} \\
\midrule
\endhead

\bottomrule
\endfoot

\texttt{MyReferenceDesignator}
    & string
    & Subsystem reference designator (e.g., \texttt{DR1}).
      Messages with non-matching designators (other than
      \texttt{ALL}) are ignored. \\

\texttt{SelfIp}
    & IP address
    & DR control interface IP address. \\

\texttt{TimeAuthority}
    & IP or hostname
    & NTP time server used for synchronization on boot
      and via the \cmd{SYN} command. \\

\texttt{SerialNumber}
    & string
    & Serial number identifying this DR instance. \\

\texttt{ArraySelect}
    & string
    & DRSU to select on startup.  May be a 0-based index or a
      DRSU barcode. \\

\texttt{MessageInPort}
    & integer
    & UDP port for receiving MCS messages. \\

\texttt{MessageOutPort}
    & integer
    & UDP port for sending MCS responses. \\

\texttt{DataInIp}
    & IP address
    & IP address to bind for receiving NDP data frames.
      Defaults to the value of \texttt{SelfIp}. \\

\texttt{DataInPort}
    & integer
    & UDP port for receiving NDP data frames. \\

\end{longtable}
